\documentclass[10pt,a4paper]{amsart}
\usepackage[margin=19mm]{geometry}
\usepackage{amsmath,amssymb,mathtools}
\usepackage[scr=rsfs,cal=euler]{mathalfa}
\usepackage{xfrac}
\usepackage[unicode,hidelinks,bookmarks=false]{hyperref}
\newcommand{\ie}{i.e.\ }
\newcommand{\cf}{cf.\ }
\newcommand{\resp}{respectively\ }
\newcommand{\Subsection}{\subsection}
\providecommand{\nobreakdash}{\nobreak\hbox{-}}
\setlength{\emergencystretch}{2em}
\multlinegap0pt
\usepackage{bm}
\usepackage{bbm}
\usepackage{dsfont}
\usepackage{xspace}
\usepackage{cancel}
\usepackage{mathtools}
\usepackage{amsthm}
\makeatletter
\@ifundefined{lemma}{\newtheorem{lemma}{Lemma}}{}
\@ifundefined{claim}{\newtheorem{claim}[lemma]{Claim}}{}
\@ifundefined{proposition}{\newtheorem{proposition}[lemma]{Proposition}}{}
\makeatother
\newcommand{\CI}{{\mathcal I}}
\newcommand{\CX}{{\mathcal X}}
\newcommand{\E}{{\mathbb E}}
\newcommand{\N}{{\mathbb N}}
\newcommand{\R}{{\mathbb R}}
\newcommand{\Z}{{\mathbb Z}}
\newcommand{\floor}[1]{{\left \lfloor #1 \right \rfloor}}
\newcommand{\norm}[1]{\left\Vert #1\right\Vert}
\newcommand{\rem}[1]{\left \{ #1 \right \}}
\title[Resolving the joint ergodicity problem for Hardy sequences]{Resolving the joint ergodicity problem for Hardy sequences}
\author{Sebasti\'an Donoso, Andreas Koutsogiannis, Borys Kuca, Wenbo Sun, Konstantinos Tsinas}
\date{Source: arXiv:2506.20459v2 (CC BY 4.0)}
\begin{document}
\maketitle

\noindent\emph{Source and licence.} Resolving the joint ergodicity problem for Hardy sequences. Version arXiv:2506.20459v2 (CC BY 4.0), distributed under the Creative Commons Attribution 4.0 International licence, as recorded in the arXiv licence metadata for this exact version and shown on the abstract page https://arxiv.org/abs/2506.20459v2. Full rights notice: licenses/arxiv-2506.20459-CC-BY-4.0.txt.\par\smallskip

\noindent\textit{Editorial scope.} Counted unit: the Proposition whose source label is \texttt{P: beta rational} in the article (source0.tex lines 2822--2824, source label \texttt{P: beta rational}), delivered together with the article's own complete original proof of it (source0.tex lines 2850--3017, one \texttt{proof} environment of 168 lines). No mathematics is written, repaired or re-derived: statement and proof are the article's own text line for line. The proof is detached from the statement in the source: the article states the result at lines 2822--2824 and prints its proof at lines 2850--3017, and the proof environment's own header names the statement, so the association is the article's own. Required context retained uncounted: none. Every definition, display and label of those lines is retained unchanged. Omitted from this package and not redistributed: 1--2821, 2825--2849, 3018--4790. Nothing inside a retained excerpt is omitted or altered: every retained body is delivered exactly as the article's lines are written, so no omission receipt of any kind accompanies this package. Citations: the retained excerpts carry no citation command at all, so no bibliography entry of the article is transcribed and no reference is redistributed. Reference adaptation: none was needed and none was made; every reference inside the delivered unit resolves inside it, so no unresolved reference or citation is produced. Printed slips kept literal: no printed word, symbol, hypothesis or number of the counted statement or of its proof was altered. Layout and class: the article's own class is \texttt{amsart} with packages including amssymb, amsfonts, hyperref, mdwlist, textcomp, enumerate, inputenc, fontenc, eucal, xcolor, geometry, comment; the wrapper is the lane's pinned \texttt{amsart} editorial wrapper at 10pt on A4, which restates the theorem-like environments the retained text needs and adds the article's own macro definitions that the retained text needs (\texttt{\textbackslash CI}, \texttt{\textbackslash CX}, \texttt{\textbackslash E}, \texttt{\textbackslash N}, \texttt{\textbackslash R}, \texttt{\textbackslash Z}, \texttt{\textbackslash floor}, \texttt{\textbackslash norm}, \texttt{\textbackslash rem}), with extra wrapper packages (bm, bbm, dsfont, xspace, cancel, mathtools). The delivered numbering is the unit's own. Assets: no retained span contains a figure environment, a graphics-inclusion command, a PDF-page inclusion, a table or a TikZ picture; no image file is copied into this package, the package declares no asset dependency and no image file is redistributed with it. Licence: version arXiv:2506.20459v2 of this article is distributed under the Creative Commons Attribution 4.0 International licence, as recorded in the arXiv licence metadata for this exact version and shown on the abstract page https://arxiv.org/abs/2506.20459v2. A full rights notice is delivered at licenses/arxiv-2506.20459-CC-BY-4.0.txt. Characters: every retained excerpt is pure ASCII, so the delivered engine, XeLaTeX with its default Latin Modern text font, reports no missing character.
\begin{proposition}\label{P: beta rational}
    Let $k,l\in\mathbb{Z}\backslash\{0\}$ and $u\in\R$. Then there exists a finite index subgroup $H\subseteq\mathbb{Z}^{2}$, which contains $(k,l)$ and depends only on $k$ and $l$, such that for any system $(X,\CX,\mu, T_{1},T_{2})$ and any sequence $a\colon\mathbb{N}\to\mathbb{R}$, the ergodicity of $(T_{1}^{\floor{ka(n)}}T_{2}^{-\floor{la(n)+u}})_{n}$ implies that $\CI(H)=\CI(T_{1}^{k}T_{2}^{-l})$.   
\end{proposition}

\begin{proof}[Proof of Proposition \ref{P: beta rational}]
Rescaling $k,l$ and $a$ if necessary, we may assume without loss of generality that $k$ is coprime to $l$. Let $I_1, \ldots, I_K$ be the finite partition of $[0,1)$ into all the intervals with endpoints from the set
    \begin{align}\label{E: endpoints}
        \Bigl\{\frac{i}{\vert k\vert},\frac{j}{\vert l\vert}, \frac{j-\{u\}}{\vert l\vert}\colon 0\leq i\leq \vert k\vert,\; 0\leq j\leq \vert l\vert\Bigr\}.
    \end{align}
% \BK{(This set is slightly different from the one that Wenbo originally wrote because of the presence of $u$. Please check if this makes sense.)} {\color{red} should be
% \begin{align}%\label{E: endpoints}
%         \Bigl\{\frac{i}{\vert k\vert},\frac{j}{\vert l\vert}, \Bigl\{\frac{j}{\vert l\vert}-u\Bigr\}\colon 0\leq i\leq \vert k\vert,\; 0\leq j\leq \vert l\vert\Bigr\}.
%     \end{align}
% } \BK{Are you sure? My computations give
% \begin{align*}
%     \floor{la(n) + u} = \floor{la(n)}+\floor{u} + 1_{\{la(n)\}+\{u\}\geq 1}.
% \end{align*}
% Now, if $\{a(n)\}\in[\frac{j}{l},\frac{j+1}{l})$, then $\{la(n)\} = l\{a(n)\}-j$ and $\floor{l a(n)}=l\floor{a(n)}+j$, and so we have
% \begin{align*}
%     \floor{la(n) + u} = l\floor{a(n)}+\floor{u} + j + 1_{l\{a(n)\}-j+\{u\}\geq 1}.
% \end{align*}
% Then the condition can be rewritten as $\{a(n)\}\geq \frac{j+1-\{u\}}{l}$.
% }
% {\color{red} I see. But then it should be $\{\frac{j-\{u\}}{\vert l\vert}\}$ since it belongs to [0,1)}
By the Bolzano-Weierstrass theorem, there exists a subsequence $(N_j)_j$ of $\mathbb{N}$ such that
the limit 
\begin{align}\label{E: weights}
    c_i =\lim_{j\to\infty}\frac{\vert\{n\in[N_j]\colon\; \{a(n)\}\in I_i\}\vert}{N_{j}}
\end{align}
    exists for all $i\in[K]$.


    % $$\lim_{j\to\infty}\frac{\vert\{n\in[N_j]\colon\; \{a(n)\}\in [x,y)\}\vert}{N_{j}}$$
    % exists for all $x,y$ in
    % \begin{align}\label{E: endpoints}
    %     \Bigl\{\frac{i}{\vert k\vert},\frac{j}{\vert l\vert}\colon 0\leq i\leq \vert k\vert, 0\leq j\leq \vert l\vert\Bigr\}.
    % \end{align}

%Denote $m:=(k,-l)$. 
From now on, assume that $f$ has zero integral and is invariant under the transformation $R:=T_1^kT_2^{-l}$. Our objective is to show that $f$ is $\CI(H)$-measurable for some fixed finite index subgroup $H\subseteq \Z^2$ {which contains $(k,l)$}.

Since $(T_{1}^{\floor{ka(n)}}T_{2}^{-\floor{la(n)+u}})_{n}$ is ergodic, we have
\begin{equation}\label{kl1}
    \lim_{j\to\infty}\E_{n\in[N_j]}T_{1}^{\floor{ka(n)}}T_{2}^{-\floor{la(n)+u}}f=0
\end{equation}
in $L^2(\mu)$.
%for any zero-integral function $f$.
%For convenience write $T_{(x,y)}:=T_{1}^{x}T_{2}^{y}$ for all $x,y\in\mathbb{Z}$.
Using the identities
\begin{gather*}
 \floor{n x} = n\floor{x} + i \quad\textrm{whenever}\quad n\in\N,\; x\in\left[\frac{i}{n}, \frac{i+1}{n}\right)\\
 \textrm{and}\quad  \floor{x + y}=\floor{x}+\floor{y}+1_{\rem{x}+\rem{y}\geq 1},
\end{gather*}
% $\floor{n x} = n\floor{x} + i$ whenever $n\in\N$ and $x\in\left[\frac{i}{n}, \frac{i+1}{n}\right)$ as well as $\floor{x + y}=\floor{x}+\floor{y}+1_{\rem{x}+\rem{y}\geq 1}$, 
we can find a function $r:[0,1)\to\Z^2$, constant on each $I_1, \ldots, I_K$ and satisfying the bound
\begin{align}\label{E: bound on r}
\norm{r}_\infty \leq  \max(|k|-1, |l|),    
\end{align}
 such that 
\begin{align*}
    T_{1}^{\floor{ka(n)}}T_{2}^{-\floor{la(n)+u}}=T_{1}^{k\floor{a(n)} + r_1(n)}T_{2}^{-l\floor{a(n)}-\floor{u}-r_2(n)} = R^{\floor{a(n)}}T_1^{r_1(n)}T_2^{-\floor{u}-r_2(n)},
\end{align*}
where $r(n)=(r_1(n), r_2(n))$.
Composing (\ref{kl1}) with $T_2^{\floor{u}}$ and using the $R$-invariance of $f$, we obtain that
\begin{equation}\label{kl2}
    0=\lim_{j\to\infty}\E_{n\in[N_{j}]}T_1^{r_1(n)}T_2^{-r_2(n)}f=\sum_{i=1}^{K}c_{i}\cdot T_1^{r_{1,i}}T_2^{-r_{2,i}} f,
\end{equation}
%for all zero-integral $R$-invariant functions $f\in L^\infty(\mu)$, 
where %$R:=T_1^kT_2^{-l}$, 
the pair $(r_{1,i}, r_{2,i})$ is the value of $r$ on $I_i$ and $c_i$ is defined in \eqref{E: weights}.

Since $k$ is coprime to $l$, there exist $w,v\in\mathbb{Z}$ with $\max(\vert w\vert, |v|)=O_{k,l}(1)$ such that $wl-vk=1$. For all $i\in[K]$, we can then write $(r_{1,i}, r_{2,i}) = d_{i}(w,v)+d'_{i}(k,-l)$ for a unique choice of $d_i, d'_i\in\Z$, which are of size $O_{k,l}(1)$ by \eqref{E: bound on r}. Denote $S:=T_1^w T_2^v$. Using the $R$-invariance of $f$ again, we can recast
$T_1^{r_{1,i}}T_2^{-r_{2,i}} f = S^{d_i} f$ for all $i\in[K]$. % and all $R$-invariant functions $f$.
So it follows from (\ref{kl2}) that 
\begin{equation}\label{kl3}
    0=\sum_{i=1}^{K}c_{i}S^{d_{i}}f. %\text{ for all zero-integral $T_{m}$-invariant function $f$.}
\end{equation}
Denote 
$$L(f):=\sum_{i=1}^{K}c_{i}S^{d_{i}+d_{\ast}}f$$
for all $f$, where $d_{\ast}=\max\{-d_{1},\dots,-d_{K},0\}$; clearly, $L(f) = 0$. 
% Then (\ref{kl3}) implies that 
% \begin{equation}\label{kl4}
%     L(f)=0 \text{ for all zero-integral $T_{m}$-invariant function $f$.}
% \end{equation}
By the definition of $d_*$ and the bounds on $d_i$, we have $0\leq d_{i}+d_{\ast} \ll_{k,l}1$, and hence
$$g(z):=\sum_{i=1}^{K}c_{i}z^{d_{i}+d_{\ast}}$$ is a real polynomial of degree $O_{k,l}(1)$. Since $c_{i}\geq 0$ and at least one $c_{i}$ is nonzero, the polynomial $g$ is nonzero. 

Thus, the study of the limit $L(f)$ can be reduced to the examination of the polynomial $g$. The point of the next few claims is that we can throw away all the roots of $g$ that are not roots of unity, and that we can combine them all together to find a finite-index subgroup $H\subseteq\Z^2$ containing $(k,-l)$ such that $f$ is $H$-invariant.
%Our goal now is to establish more properties of the polynomial $g$ that will be crucial in completing the proof of the proposition.
    
\begin{claim}\label{Claim 1}
% Let $f$ be a nonzero $T_{m}$-invariant function with
Suppose that $f\neq 0$ and $Sf=\lambda f$ for some $\lambda\in\mathbb{C}$.
Then $\lambda$ must be a root~of~$g$ and a root of unity. 
\end{claim}
%{\color{red} in the claims, it is better to say that it works for all zero-integral $R$-invariant function, not just for a fixed $f$. (or another way is to say clearly that \eqref{kl3} holds for all zero-integral $R$-invariant function)}

Indeed, assume that $\lambda$ is not a root of unity. Then for all $n\in \N$, $$\int f^n\; d\mu = \int S(f^n)\; d\mu= \int (Sf)^n\; d\mu=\lambda^n\int f^n\; d\mu,$$ so $f^n$ has zero integral.  
Since $f^n$ is $R$-invariant by assumption, we get $L(f^n)=0$ by \eqref{kl3}. Because $f^n\neq 0$, we further deduce that $\lambda\neq 0$ and $g(\lambda^n)=0$. By the fundamental theorem of algebra, the polynomial $g$ has only finitely many roots, and so the only way in which $g(\lambda^n)=0$ can happen for all $n\in\N$ is if $\lambda$ is a root of unity.

\

As an immediate corollary, we get the following.

\begin{claim}\label{Claim 2}
   Suppose that $\lambda\in\mathbb{C}$ is not a root of unity. If $(S-\lambda)f=0$, then $f=0$.
\end{claim}






 %We are now ready to complete the proof the proposition. Let $f$ be any zero-integral $T_{m}$-invariant function.
 Next, we move on to show that we can remove all the roots of $g$ that are not roots of unity. Since $g\neq 0$, we may factorize $g$ as 
 $$g(z)=c\prod_{j=1}^{M}(z-z_{j})$$
 for some $c\neq 0, M=O_{k,l}(1)$ and $z_{1},\dots,z_{M}\in\mathbb{C}$. Then it follows from (\ref{kl3}) that
 $$L(f)=c\prod_{j=1}^{M}(S-z_{j})f=0.$$

If $M=0$, then since $c\neq 0$, we must have $f=0$, in which case we are done. So we assume that $M>0$. Rearranging the roots, we can assume that  $z_{1},\dots,z_{M'}$ are roots of unity while $z_{M'+1},\dots,z_{M}$ are not. By repeatedly applying Claim \ref{Claim 2} to $\prod_{j=1}^{M'+t}(S-z_{j})f$ in place of $f$, we deduce that
$$\prod_{j=1}^{M'+t}(S-z_{j})f=0$$
for all $0\leq t\leq M-M'-1$. Hence
 $$\prod_{j=1}^{M'}(S-z_{j})f=0.$$
 %Since $Sf=0$ implies that $f=0$, we may further assume that all of $z_{1},\dots,z_{M'}$ are  roots of unity. 
 
 Assume that $z_j$ is a primitive $q_j$-th root of unity for all $1\leq j\leq M'$.
Then  
\begin{equation}\label{kl5}
    \prod_{j=1}^{M'}(S^{q_{j}}-1)f=0
\end{equation}
since $z-z_{j}$ divides $z^{q_{j}}-1$.

The following claim allows us to bound the size of $q_j$'s, and hence the index of the subgroup $H$ later on.
\begin{claim}\label{Claim 3}
    We have $q_{j}=O_{k,l}(1)$.
\end{claim}
 
% \BK{(This is essentially what Wenbo wrote modulo the proof of $\lim\limits_{t\to\infty}\phi(t)=\infty$, which I commented out. However, I am uncertain about what's going on in the argument below. Why does $\phi$ matter? Are you trying to show that the polynomial $\prod_{j=1}^{M'}(z^{q_{j}}-1)$ divides $g$, and hence it must have a degree at most $O_{k,l}(1)$? I don't see why this would be true given that a priori, $g$ may have arbitrary nonnegative real coefficients.)}{\color{red} Since  $z_j$ is a primitive $q_j$-th root of unity, its minimal polynomial $h$ is of degree $\phi(q_{j})$. Since $z_{j}$ is also a root of $g$, we must have that $h\vert g$ and so $\phi(q_{j})=\deg(h)\leq \deg(g)=O_{k,l}(1)$ unless $g\equiv 0$. Since $\lim\limits_{t\to\infty}\phi(t)=\infty$, $\phi(q_{j})=O_{k,l}(1)$ also implies that $q_{j}=O_{k,l}(1)$.} 
 To see this, recall that since  $z_j$ is a primitive $q_j$-th root of unity, its minimal polynomial $h$ is of degree $\phi(q_{j})$, where $\phi$ is the Euler totient function. Since $z_{j}$ is also a root of $g$, the polynomial $h$ must divide $g$, and so $\phi(q_{j})=\deg(h)\leq \deg(g)=O_{k,l}(1)$. Since $\lim\limits_{t\to\infty}\phi(t)=\infty$, we have $\phi(q_{j})=O_{k,l}(1)$ only if $q_{j}=O_{k,l}(1)$.
 
 
 %let $\phi(t)$ be the Euler totient function, i.e. the number of $j\in[t]$ coprime to $t$. It is known that the minimal polynomial for the $t$-th root of unity is of degree $\phi(t)$, and that $\lim\limits_{t\to\infty}\phi(t)=\infty$. The claim follows since  $z_{j}$ is both a root of unity and a root of a polynomial of degree $O_{k,l}(1)$.

%  Indeed, let $t=p_{1}^{m_{1}}\dots p_{r}^{m_{r}}$ be the prime factorization of $t$ with $p_{1}<\dots<p_{r}$. Then 
%  $$t-\phi(t)=(\sum_{i=1}^{r}\frac{1}{p_{i}}-\sum_{1\leq i_{1}<i_{2}\leq r}\frac{1}{p_{i_{1}}p_{i_{2}}}+\dots+(-1)^{r}\frac{1}{p_{1}\dots p_{r}})t$$
% and so
%  $$\phi(t)=t\prod_{i=1}^{r}(1-\frac{1}{p_{i}})=\prod_{i=1}^{r}(p_{i}-1)p_{i}^{m_{i}-1}.$$
%  Since $(p_{i}-1)p_{i}^{m_{i}-1}\geq p_{i}^{m_{i}/2}$ if $m_{i}\geq 2$ or if $m_{i}=1$ and $p_{i}\geq 3$. We have that $\phi(t)\geq \frac{1}{100}\sqrt{t}$. This completes the proof of Claim 3.
 
% \

The next claim enables us to combine the contributions of all the roots of unity $z_1, \ldots, z_{M'}$.

\begin{claim}\label{Claim 4}
Let $m,n\in\mathbb{N}$. If $(S^{m}-1)(S^{n}-1)f=0$, then $(S^{2mn}-1)f=0$.    
\end{claim}

To see this, note that $(S^{m}-1)(S^{n}-1)f=0$ implies that $(S^{mn}-1)^{2}f=0$ since $z^{n}-1$ and $z^{m}-1$ both divide $z^{mn}-1$. Then 
$(S^{2mn}+1)f=2S^{mn}f.$
     So $$4\int \vert f\vert^{2}=4\int \vert S^{mn}f\vert^{2}=\int \vert (S^{2mn}+1)f\vert^{2}=2\int \vert f\vert^{2}+\langle S^{2mn}f,f\rangle+\langle f,S^{2mn}f\rangle.$$
     By the Cauchy-Schwarz inequality, $\langle S^{2mn}f,f\rangle\leq \vert f\vert^{2}$ and the equality holds only if $S^{2mn}f=f$. Hence $(S^{2mn}-1)f=0$, and the claim follows. 

     \

     %\BK{(I changed the quantities below a bit. Note that the choice of $q$ depends only on $g$ (and hence on $a,k,l,u$) and not on $f$)} {\color{red} If we take $q=2^{M'-1}(O_{k,l}(1)!)^{M'}$, then $q$ depends only on $k,l$}
     Recall that $M'\leq M=O_{k,l}(1)$ and $q_{j}=O_{k,l}(1)$, the latter of which follows from Claim \ref{Claim 3}. A repeated application of Claim \ref{Claim 4} and (\ref{kl5}) allows us to find $W=O_{k,l}(1)$ and $q = 2^{M'-1} q_1 \cdots q_{M'}\leq W$ such that $S^{q}f=f$.
     %there exist $q, W=O_{k,l}(1)\in \N$ with $q\leq W$ such that 
     Let $H = \langle W! (w,v), (k,-l)\rangle$: this is a finite index subgroup of $\mathbb{Z}^{2}$ that depends only on $k$ and $l$. 
     Since $q|W!$, the fact that $f$ is invariant under both $S^{q}$ and $R$ implies that $f$ is $H$-invariant. This means that $\CI(R)=\CI(H)$. 
     %\BK{(In our statement, we have that $H$ can be made independent of $a$. While not particularly important, we should give reason why this is the case. Presumably this would follow if we know that $\Z^2$ admits $O_q(1)$ distinct index-$q$ subgroups - is this true?)}
     %Since $M'=O_{k,l}(1)$ and since $q_{j}=O_{k,l}(1)$ by Claim 3, it follows from Claim 4 and (\ref{kl5}) that there exists $W=O_{k,l}(1)\in \N$ and $q\in \mathbb{N}_{+}$ with $q\leq W$ such that $S^{q}f=f$. Let $H$ be the $\mathbb{Z}$-span of $W!n$ and $m$, which is certainly a finite index subgroup of $\mathbb{Z}^{2}$ depending only on $k$ and $l$. The fact that $f$ is $S^{q}$ and $T_{m}$ invariant implies that $f$ is $H$ invariant. This means that $I(T_{m})=I(H)$. 
     \end{proof}

\end{document}
